%======================================================================
% HABILIDAD: Introducción General
% ID: introduccion
% Descripción: Propósito, alcance y estructura del documento
%======================================================================

\capitulo{Introducción}

\section{Propósito del Documento}

Este manual tiene como objetivo principal proporcionar una guía completa y detallada sobre el uso, configuración y administración del sistema. El documento está diseñado para servir como referencia tanto para usuarios finales como para administradores técnicos.

La información contenida en este manual permitirá al lector:
\begin{itemize}
    \item Comprender los conceptos fundamentales del sistema
    \item Realizar la instalación y configuración inicial
    \item Utilizar todas las funcionalidades disponibles
    \item Resolver problemas comunes de manera autónoma
    \item Implementar las mejores prácticas de seguridad
\end{itemize}

\section{Audiencia Objetivo}

Este manual está dirigido a las siguientes audiencias:

\begin{description}
    \item[Usuarios Finales] Personas que utilizarán el sistema para sus actividades diarias. Se asumen conocimientos básicos de computación.
    
    \item[Administradores] Responsables de la configuración, mantenimiento y soporte del sistema. Se requieren conocimientos técnicos intermedios.
    
    \item[Desarrolladores] Profesionales que integrarán el sistema con otras aplicaciones. Se requieren conocimientos avanzados de programación.
\end{description}

\section{Estructura del Manual}

El contenido de este manual está organizado de la siguiente manera:

\begin{enumerate}
    \item \textbf{Introducción} - Presenta el propósito y alcance del documento
    \item \textbf{Instalación} - Guía paso a paso para la instalación del sistema
    \item \textbf{Configuración} - Opciones de personalización y ajustes
    \item \textbf{Panel de Control} - Descripción de la interfaz principal
    \item \textbf{Gestión} - Administración de recursos y usuarios
    \item \textbf{Seguridad} - Políticas y prácticas de seguridad
    \item \textbf{Anexos} - Referencias técnicas adicionales
\end{enumerate}

\section{Convenciones Utilizadas}

A lo largo de este documento se utilizan las siguientes convenciones tipográficas:

\begin{itemize}
    \item \textbf{Texto en negrita}: Indica opciones de menú, botones o elementos importantes
    \item \texttt{Texto monoespaciado}: Representa código, comandos o nombres de archivos
    \item \textit{Texto en cursiva}: Denota términos nuevos o énfasis especial
    \item \faIcon{info-circle} Indica información adicional o consejos útiles
    \item \faIcon{exclamation-triangle} Señala advertencias o precauciones importantes
\end{itemize}

\notacientifica{Las versiones de software mencionadas en este manual corresponden a la versión {{VERSION}} del sistema. Para versiones anteriores o posteriores, algunas funcionalidades podrían variar.}

\section{Requisitos Previos}

Antes de comenzar con la instalación y uso del sistema, se recomienda contar con:

\begin{itemize}
    \item Conocimientos básicos de sistemas operativos
    \item Acceso a una cuenta de usuario con privilegios administrativos
    \item Conexión estable a Internet (para descargas y actualizaciones)
    \item Los recursos de hardware mínimos especificados en la sección de instalación
\end{itemize}
